\chapter{Conclusions and Future Work}
\label{ch:conclusions}

The reliability of an autonomous vehicle depends on how accurately it can estimate the state of
the vehicles around it, and on how this estimate copes with noisy sensors and with dynamics
that the models describe only approximately. This thesis has addressed this problem in the
context of autonomous racing, a setting in which high speeds, abrupt maneuvers and close
interactions between vehicles make the limitations of the estimation clearly evident, by
working on the target tracking module of the Dallara AV-24 and EAV-25 and evaluating the
proposed methods on experimental racing data.

The analysis of the measurement noise has shown that the errors of the detections cannot in
general be described as stationary, white Gaussian noise. Once the contributions related to
the reference and to the calibration are removed, the errors of several sources still contain
a slowly varying component, which has been modeled as a Gauss-Markov process and separated
from a white one, while their dependence on the relative geometry has proved not to be
reproducible across sessions, preventing a general deterministic compensation. These results
indicate that a correct description of the measurement noise requires its temporal structure
to be taken into account, beyond its first two moments. On this basis, a covariance matching
technique has been adopted to estimate the measurement noise covariance of the LiDARs
recursively and in real time, and it has been shown to describe the white component of the
errors more closely than a static covariance, while the slowly varying bias, shared by the
estimate of the state, cannot be recovered from the innovations alone.

The degradation of the estimate during the maneuvers has been addressed with two complementary
approaches. The feed-forward adaptation of the process noise has proved to be correct in
principle, since it makes the filter more reactive when a maneuver is detected and preserves
the nominal behavior elsewhere, but it has also shown the limits of any reactive strategy,
whose effectiveness depends on the delay and on the noise of the signal used for the
detection. The RTS smoother, instead, has shown that exploiting the measurements acquired
after each instant considerably improves the estimate of the speed, both with and without a
direct measurement of it, at the price of a state referred to the past, which is acceptable
whenever the downstream modules, such as the gap control of the planning module, do not
strictly require the current state. The combination of the two approaches has provided the
most accurate estimate among the configurations considered, confirming that improving the
causal estimate remains beneficial even when it is subsequently refined by the smoother.

More generally, these results show that a careful modeling of the uncertainties of a tracking
filter, both on the measurements and on the motion model, can considerably improve the
estimate without modifying the architecture of the system, and that the trade-off between
latency and accuracy can be exploited whenever the application tolerates a delayed estimate.
Both aspects are not specific to racing, and they apply equally to the tracking of the
vehicles in urban traffic, where heterogeneous sensors and sudden changes of the dynamics of
the other road users pose the same problems.

\section{Future Work}
\label{sec:future_work}

As anticipated in Section~\ref{sec:ch4_future}, a first extension concerns the measurement
noise of the camera and of the RaDAR. Their innovations contain, besides their own white
error, the difference between their bias and the one of the LiDARs, which varies slowly and
prevents the covariance matching from being applied to them directly. Unlike the absolute
bias, this relative bias is observable, and it could be estimated online by augmenting the
state of the filter with a slowly varying bias for each sensor, modeled as a Gauss-Markov
process. Once compensated, the innovations of the camera and of the RaDAR would contain only
their white component, and the adaptive estimation of the measurement noise covariance could
be extended to all the detection sources, improving the relative weighting of their
measurements in the filter.

A second extension concerns the feed-forward adaptation of the process noise, whose
performance is limited by the trade-off between delay and noise of the jerk estimate obtained
by differentiating and filtering the acceleration. A constant-jerk motion model could be
implemented, extending the constant-acceleration model with the jerk as an additional state of
the filter, so that the jerk would be estimated directly by the filter instead of being
derived from its output. It would then be possible to evaluate whether this formulation
resolves the trade-off between delay and noise, providing an earlier and more reliable
detection of the maneuvers, in terms of activation delay, false activations and estimation
error.

Finally, the proposed methods should be validated on a larger set of sessions, circuits and
operating conditions, and their effect on the downstream modules, above all on the planning
and on the gap control, should be evaluated at the level of the whole autonomous driving
system.
